\chapter{Metodologia} \label{chap:metodologia}

\section{Definição da estrutura de tópicos \mqtt{}}

A primeira etapa do trabalho consistirá em definir uma estrutura de tópicos no \broker{} \mqtt{}, onde as mensagens geradas pelos eventos de descoberta, conexão e desconexão serão publicadas. A estrutura de tópicos deve permitir flexibilidade para as aplicações indicarem interesse nos eventos.

\section{Implementação no \middleware{} \mhubcddl{}}

O trabalho seguirá com a implementação de um novo módulo no \mhubcddl{} que gera os eventos de descoberta, conexão e desconexão, e os publica em um \broker{} \mqtt{} seguindo a estrutura de tópicos defina ateriormente.

Também será desenvovido uma interface que permite que os desenvolvedores indiquem que suas aplicações possuem interesse em determinados eventos.

\section{Teste e validação}

Serão avaliados parâmetros de acurácia e performance da solução. A acurácia representa a quantidade de notificações efetivamente entregues à aplicação, e o performance se refere a velocidade com que a entrega das mensagens é feita.

Para realizar as medições mencionadas, será utilizado um caso de uso simulado com um \dataset{}, onde um indivíduo se locomove em sua residencia enquanto seu \smartphone{} descobre e se conecta com os \smartobjs{} localizados na casa. As notificações geradas neste caso de uso serão então analizadas para determinar a porcentagem de perdas e o tempo em que elas levam até chegaram na camada de aplicação.

Como resultado da investigação preliminar, destaca-se o \dataset{} disponibilizado por \citeonline{byrne2018residential} como potencialmente relevante para os testes propostos.

